\documentclass{article}

% if you need to pass options to natbib, use, e.g.:
%     \PassOptionsToPackage{numbers, compress}{natbib}
% before loading neurips_2024



% ready for submission
%\usepackage{neurips_2024}


% to compile a preprint version, e.g., for submission to arXiv, add add the
% [preprint] option:
\usepackage[preprint]{neurips_2024}


% to compile a camera-ready version, add the [final] option, e.g.:
%\usepackage[final]{neurips_2024}


% to avoid loading the natbib package, add option nonatbib:
%    \usepackage[nonatbib]{neurips_2024}


\usepackage[utf8]{inputenc} % allow utf-8 input
\usepackage[T1]{fontenc}    % use 8-bit T1 fonts
\usepackage{hyperref}       % hyperlinks
\usepackage{url}            % simple URL typesetting
\usepackage{booktabs}       % professional-quality tables
\usepackage{amsfonts}       % blackboard math symbols
\usepackage{nicefrac}       % compact symbols for 1/2, etc.
\usepackage{microtype}      % microtypography
\usepackage{xcolor}         % colors
\usepackage{amsmath}
\usepackage{graphicx}


\title{TBD}


% The \author macro works with any number of authors. There are two commands
% used to separate the names and addresses of multiple authors: \And and \AND.
%
% Using \And between authors leaves it to LaTeX to determine where to break the
% lines. Using \AND forces a line break at that point. So, if LaTeX puts 3 of 4
% authors names on the first line, and the last on the second line, try using
% \AND instead of \And before the third author name.


\author{%
  Simon ~Danninger\\
  Institute for Machine Learning\\
  Johannes Kepler University\\
  Linz, Austria \\
  \texttt{simon.danninger64@gmail.com} \\
}


\begin{document}


\maketitle


\begin{abstract}

\end{abstract}


\section{Introduction}




\section{Background \& Related Work}


\section{Datasets and Experimental Setup}
We perform our experiments on four datasets from The Well and PDEBench. Our aim is to cover diverse time dependent physics in one and two dimensions and we restrict ourselves to cartesian domains. An Overview of the datasets is given in Table \ref{tab:datasets}.


\begin{table}[h]
  \centering

  \begin{tabular}{llccccc}
    \toprule
    Dataset & Source & $N_d$ & Resolution & Fields & $N_\mathrm{steps}$ & $N_\mathrm{samples}$ \\
    \midrule
    \texttt{Active Matter}                  & The Well & 2 & $256\times256$ & 11 & 81  & 360   \\
    \texttt{Turbulent Radiative Layer} & The Well & 2 & $128\times384$ & 4  & 101 & 90    \\
    \texttt{Shallow Water Equation}                   & PDEBench & 2 & $128\times128$ & 1  & 101 & 1000  \\
    \texttt{Burgers'} & PDEBench & 1 & $256$          & 1  & 41  & 10000 \\
    \bottomrule
  \end{tabular}
  \vspace{0.4em}
    \caption{Overview of datasets. $N_d$ denotes the spatial dimensionality, \emph{Resolution} the per-snapshot grid, \emph{Fields} the number of physical channels, $N_\mathrm{steps}$ the number of time-steps per trajectory, and $N_\mathrm{samples}$ the number of trajectories.}
  \label{tab:datasets}
\end{table}

\subsection{Active Matter}
The Active Matter (AM) dataset from The Well \citep{ohana2024well, maddu2024learning} simulates a dense suspension of active, rod-like particles in a viscous fluid, where particle activity and alignment drive chaotic large-scale flows.

The dynamics follow a kinetic theory in which the particle distribution function $\Psi(\mathbf{x}, \mathbf{p}, t)$, defined over position $\mathbf{x}$ and orientation $\mathbf{p}$ with $|\mathbf{p}| = 1$, is advected by a Stokes flow:
\begin{align}
  \frac{\partial \Psi}{\partial t} + \nabla_{\mathbf{x}} \cdot (\dot{\mathbf{x}}\,\Psi) + \nabla_{\mathbf{p}} \cdot (\dot{\mathbf{p}}\,\Psi) &= 0, \\
  \dot{\mathbf{x}} &= \mathbf{u} - d_T\, \nabla_{\mathbf{x}} \log \Psi, \\
  \dot{\mathbf{p}} &= (\mathbf{I} - \mathbf{p}\mathbf{p}) \cdot (\nabla \mathbf{u} + 2\zeta \mathsf{D}) \cdot \mathbf{p} - d_R\, \nabla_{\mathbf{p}} \log \Psi, \\
  -\Delta \mathbf{u} + \nabla P &= \nabla \cdot \boldsymbol{\Sigma}, \qquad \nabla \cdot \mathbf{u} = 0, \\
  \boldsymbol{\Sigma} &= \alpha \mathsf{D} + \beta\, \mathsf{S}\!:\!\mathsf{E} - 2\zeta\beta\,(\mathsf{D}\cdot\mathsf{D} - \mathsf{S}\!:\!\mathsf{D}).
\end{align}
Here $\mathbf{u}$ is the fluid velocity, $P$ the pressure, and $\boldsymbol{\Sigma}$ the active stress tensor. The orientation moments $\mathsf{D} = \langle \mathbf{p}\mathbf{p} \rangle$ and $\mathsf{S} = \langle \mathbf{p}\mathbf{p}\mathbf{p}\mathbf{p} \rangle$ are the second and fourth moments of $\Psi$, $\mathsf{E} = \tfrac{1}{2}(\nabla \mathbf{u} + \nabla \mathbf{u}^{\top})$ is the rate-of-strain tensor, and $\mathbf{I}$ is the identity. The constants $d_T$ and $d_R$ are translational and rotational diffusion coefficients. The simulation records the low-order moments of $\Psi$, namely the concentration $c = \langle 1 \rangle$, the velocity $\mathbf{u}$, the orientation tensor $\mathsf{D}$, and the strain-rate tensor $\mathsf{E}$, which together form the 11 channels listed in Table~\ref{tab:datasets}.

The domain is a periodic square of side $L = 10$ on a $256 \times 256$ grid, and each trajectory holds 81 snapshots. Two parameters are varied: the active dipole strength $\alpha \in \{-1, -2, -3, -4, -5\}$ and the alignment strength $\zeta \in \{1, 3, 5, 7, 9, 11, 13, 15, 17\}$, while the packing parameter is fixed at $\beta = 0.8$. The resulting 45 parameter pairs, each simulated from several random initial conditions, make up the 360 trajectories of the dataset.

\subsection{Turbulent Radiative Layer}
The Turbulent Radiative Layer (TRL) dataset from The Well \citep{ohana2024well, fielding2020multiphase} simulates the turbulent mixing layer that forms between cold dense gas and hot dilute gas, where radiative cooling competes with turbulent mixing.

The dynamics are governed by the compressible Euler equations with a radiative cooling source term:
\begin{align}
  \frac{\partial \rho}{\partial t} + \nabla \cdot (\rho \mathbf{v}) &= 0, \\
  \frac{\partial (\rho \mathbf{v})}{\partial t} + \nabla \cdot (\rho\, \mathbf{v} \otimes \mathbf{v} + P\mathbf{I}) &= 0, \\
  \frac{\partial E}{\partial t} + \nabla \cdot \big((E + P)\mathbf{v}\big) &= -\frac{E}{t_{\mathrm{cool}}}, \qquad E = \frac{P}{\gamma - 1},
\end{align}
where $\rho$ is the gas density, $\mathbf{v}$ the velocity, $P$ the pressure, and $E$ the energy density with adiabatic index $\gamma = 5/3$. The source term removes thermal energy on a cooling timescale $t_{\mathrm{cool}}$ that is shortest at intermediate temperatures, so freshly mixed gas cools rapidly and rejoins the cold phase.

The domain is Cartesian with $x \in [-0.5, 0.5]$ and $y \in [-1, 2]$, discretised on a $128 \times 384$ grid with 101 snapshots per trajectory. The boundary conditions are periodic in $x$ and zero-gradient (outflow) in $y$. The cooling time is the varied parameter, $t_{\mathrm{cool}} \in \{0.03, 0.06, 0.1, 0.18, 0.32, 0.56, 1.0, 1.78, 3.16\}$, and each of these 9 values is simulated with 10 random seeds, giving 90 trajectories. The recorded fields are the density, the pressure, and the two velocity components (4 channels).

\subsection{Shallow Water Equation}
The Shallow Water Equation (SWE) dataset from PDEBench \citep{takamoto2022pdebench} simulates a two-dimensional radial dam break, in which a raised circular column of water collapses and spreads outward as shallow-water waves.

The dynamics follow the 2D shallow water equations:
\begin{align}
  \frac{\partial h}{\partial t} + \frac{\partial (hu)}{\partial x} + \frac{\partial (hv)}{\partial y} &= 0, \\
  \frac{\partial (hu)}{\partial t} + \frac{\partial}{\partial x}\!\left(u^2 h + \tfrac{1}{2} g_r h^2\right) &= -g_r h\, \frac{\partial b}{\partial x}, \\
  \frac{\partial (hv)}{\partial t} + \frac{\partial}{\partial y}\!\left(v^2 h + \tfrac{1}{2} g_r h^2\right) &= -g_r h\, \frac{\partial b}{\partial y},
\end{align}
where $h$ is the water depth, $(u, v)$ the horizontal velocities, $b$ the bathymetry, and $g_r$ the gravitational acceleration. The bottom is flat, so $b = 0$ and the source terms vanish, and $g_r = 1$. Each simulation starts from a radial dam break with $h(0, x, y) = 2.0$ inside a central circle of radius $r$ and $h = 1.0$ outside, where the dam radius is sampled per trajectory from $r \sim \mathcal{U}(0.3, 0.7)$. The domain is $[-2.5, 2.5]^2$ with reflective (Neumann) walls, discretised on a $128 \times 128$ grid with 101 snapshots. PDEBench stores only the water depth $h$, so the dataset has a single field, and we use all 1000 trajectories.

\subsection{Burgers'}
The Burgers' dataset from PDEBench \citep{takamoto2022pdebench} simulates the one-dimensional viscous Burgers equation, a canonical model in which nonlinear advection is balanced by diffusion and develops sharp fronts at low viscosity.

The governing equation is
\begin{equation}
  \frac{\partial u}{\partial t} + \frac{\partial}{\partial x}\!\left(\frac{u^2}{2}\right) = \frac{\nu}{\pi}\,\frac{\partial^2 u}{\partial x^2}, \qquad x \in (0, 1),\; t \in (0, 2],
\end{equation}
where $u(t, x)$ is the velocity field and $\nu$ is the diffusion coefficient. We use the low-viscosity setting $\nu = 0.001$, which produces shock-like fronts and makes the dynamics hard to roll out. The boundary conditions are periodic, and each of the 10{,}000 initial conditions is a superposition of sinusoidal modes with random amplitudes and phases.

PDEBench provides this dataset at a spatial resolution of 1024 points with 201 time steps. Following PDEBench, we subsample it by a factor of 4 in space (1024 to 256 points) and by a factor of 5 in time (201 to 41 steps). We adopt the same reduction so that our results remain directly comparable to the FNO and U-Net baselines reported by PDEBench, which are trained and evaluated on the identically reduced data.

\section{Methodology}


\section{Results}


\section{Discussion}

\section{Conclusion}

{
\small

\bibliographystyle{abbrv}
\bibliography{references}
}


%%%%%%%%%%%%%%%%%%%%%%%%%%%%%%%%%%%%%%%%%%%%%%%%%%%%%%%%%%%%
\newpage

\appendix

\section{Appendix A}




%%%%%%%%%%%%%%%%%%%%%%%%%%%%%%%%%%%%%%%%%%%%%%%%%%%%%%%%%%%%




\end{document}